% Estilo visual do relatório técnico (incluído pelo pandoc com -H)

% ---------- Fontes ----------
\usepackage{fontspec}
\setmainfont{Lato}[
  UprightFont = * Regular,
  BoldFont = * Bold,
  ItalicFont = * Italic,
  BoldItalicFont = * Bold Italic,
  Scale = 0.98
]
\setsansfont{Lato}
\setmonofont{Ubuntu Mono}[Scale = 1.0]
\newfontfamily\titlefont{Lato}[UprightFont = * Black, BoldFont = * Black]
\newfontfamily\headfont{Lato}[UprightFont = * Bold, BoldFont = * Black]
\newfontfamily\lightfont{Lato}[UprightFont = * Light]
\usepackage{microtype}
\linespread{1.12}

% ---------- Cores ----------
\definecolor{primaria}{HTML}{0F4C5C}
\definecolor{acento}{HTML}{E36414}
\definecolor{fundo}{HTML}{EEF4F5}
\definecolor{texto}{HTML}{22313A}
\definecolor{cinza}{HTML}{6B7B83}
\definecolor{link}{HTML}{1F7A8C}
\color{texto}

% ---------- Página ----------
\usepackage{parskip}
\setlength{\parskip}{0.55em}
\usepackage{lastpage}
\usepackage{fancyhdr}
\pagestyle{fancy}
\fancyhf{}
\fancyhead[L]{\small\color{cinza}\lightfont Relatório Técnico · Tech Challenge Fase 1}
\fancyhead[R]{\small\color{cinza}\lightfont Apoio ao diagnóstico com Machine Learning}
\fancyfoot[R]{\small\color{primaria}\headfont\thepage\color{cinza}\lightfont\ / \pageref*{LastPage}}
\renewcommand{\headrulewidth}{0.4pt}
\renewcommand{\headrule}{\color{fundo!60!primaria}\hrule height\headrulewidth}
\fancypagestyle{plain}{\fancyhf{}\renewcommand{\headrulewidth}{0pt}
  \fancyfoot[R]{\small\color{primaria}\headfont\thepage\color{cinza}\lightfont\ / \pageref*{LastPage}}}

% ---------- Títulos de seção ----------
\usepackage{titlesec}
\titleformat{\section}
  {\headfont\LARGE\color{primaria}}{}{0em}{}
  [\vspace{2pt}{\color{acento}\titlerule[1.2pt]}]
\titlespacing*{\section}{0pt}{1.6em}{0.9em}
\titleformat{\subsection}
  {\headfont\large\color{primaria}}{}{0em}{}
\titlespacing*{\subsection}{0pt}{1.2em}{0.35em}
\titleformat{\subsubsection}
  {\headfont\normalsize\color{acento}}{}{0em}{}

% ---------- Listas ----------
\usepackage{enumitem}
\setlist{itemsep=0.25em, topsep=0.3em, parsep=0pt}
\setlist[itemize,1]{label={\color{acento}\small$\blacktriangleright$}}
\setlist[enumerate,1]{label={\color{primaria}\headfont\arabic*.}}

% ---------- Figuras e legendas ----------
\usepackage{float}
\floatplacement{figure}{H}
\usepackage[font={small,color=cinza}, labelfont={bf,color=primaria}, labelsep=period, skip=6pt]{caption}
\captionsetup[figure]{name=Figura}
\captionsetup[table]{name=Tabela}

% ---------- Tabelas ----------
\usepackage{colortbl}
\arrayrulecolor{primaria!70}
\AtBeginEnvironment{longtable}{\small\rowcolors{2}{fundo}{white}}
\setlength{\aboverulesep}{0pt}
\setlength{\belowrulesep}{0pt}
\renewcommand{\arraystretch}{1.35}

% ---------- Código ----------
\usepackage{tcolorbox}
\tcbuselibrary{skins,breakable}
\let\oldShaded\Shaded
\let\endoldShaded\endShaded
\renewenvironment{Shaded}
  {\begin{tcolorbox}[enhanced, breakable, colback=fundo, colframe=fundo, boxrule=0pt,
     arc=3pt, left=8pt, right=8pt, top=4pt, bottom=4pt, borderline west={3pt}{0pt}{primaria}]}
  {\end{tcolorbox}}

% ---------- Caixa de destaque (blockquotes do Markdown) ----------
\newtcolorbox{destaque}{enhanced, breakable, colback=acento!7, colframe=acento!7, boxrule=0pt,
  arc=3pt, left=10pt, right=10pt, top=6pt, bottom=6pt, before skip=0.8em, after skip=0.8em,
  borderline west={3pt}{0pt}{acento}}

% ---------- Capa e sumário (chamados pelo filtro Lua) ----------
\usepackage{tikz}
% \capa{título}{subtítulo}{data}{nomes separados por \\}
\newcommand{\capa}[4]{%
  \pagenumbering{gobble}
  \thispagestyle{empty}
  \begin{tikzpicture}[remember picture, overlay]
    \fill[primaria] (current page.north west) rectangle ([yshift=-0.55\paperheight]current page.north east);
    \fill[acento] ([yshift=-0.55\paperheight]current page.north west)
      rectangle ([yshift=-0.55\paperheight-6pt]current page.north east);
    \node[anchor=south west, text width=0.75\paperwidth, align=left, inner sep=0pt]
      at ([xshift=2.6cm, yshift=-0.55\paperheight+2.4cm]current page.north west) {%
      {\color{acento}\headfont\normalsize\addfontfeatures{LetterSpace=14}TECH CHALLENGE · FASE 1}\par\vspace{14pt}
      {\color{white}\titlefont\fontsize{40}{46}\selectfont #1\par}\vspace{12pt}
      {\color{white!85!primaria}\lightfont\fontsize{17}{22}\selectfont #2\par}};
    \node[anchor=north west, text width=0.75\paperwidth, align=left, inner sep=0pt]
      at ([xshift=2.6cm, yshift=-0.55\paperheight-2.2cm]current page.north west) {%
      {\color{acento}\headfont\small\addfontfeatures{LetterSpace=14}INTEGRANTES}\par\vspace{10pt}
      {\color{texto}\large\linespread{1.5}\selectfont #4\par}};
    \node[anchor=south west, inner sep=0pt]
      at ([xshift=2.6cm, yshift=2.4cm]current page.south west) {%
      {\color{cinza}\lightfont\large Pós-Tech \textbar{} #3}};
  \end{tikzpicture}
  \clearpage}

\usepackage{titletoc}
\titlecontents{section}[0pt]
  {\addvspace{8pt}\headfont\color{primaria}}{}{}
  {\hfill\contentspage}
\titlecontents{subsection}[1.5em]
  {\addvspace{2pt}\color{texto}}{}{}
  {\color{cinza!60}\titlerule*[0.7pc]{.}\color{cinza}\contentspage}
\newcommand{\sumario}{%
  \thispagestyle{empty}
  \renewcommand{\contentsname}{Sumário}
  \tableofcontents
  \clearpage
  \pagenumbering{arabic}}

% ---------- Links ----------
\AtBeginDocument{\hypersetup{colorlinks=true, linkcolor=link, urlcolor=link, citecolor=link}}
